\documentclass[12pt]{article} 
\renewcommand{\thempfootnote}{\arabic{mpfootnote}}

\usepackage[utf8]{inputenc}

\usepackage{multicol}
\setlength{\columnsep}{1cm}

\usepackage{abc}
\usepackage{graphicx}

\usepackage[margin=0.5in]{geometry} 
\geometry{letterpaper}

\usepackage{fancyhdr} 
\pagestyle{fancy} % options: empty , plain , fancy
\renewcommand{\headrulewidth}{0pt} % customise the layout...
\lhead{}\chead{}\rhead{}
\lfoot{}\cfoot{\thepage}\rfoot{}

\begin{document}

\thispagestyle{empty}

\begin{abc}[name=forkeddeer,options={-O= +c -F single.fmt}]
X:1
T:Forked Deer
M:C|
L:1/8
K:D
(3ABc|"D"defg a2 f2|"G"gfeg fedA|"D"defg a2 f2|"A"gfed cABc |
"D"defg a2 f2|"G"gfeg fedc|"D"dAFD "G"GBAG|1"A"FDEF "D"D2:|2"A"FDEF "D"D4||
|:"A"A2 A2 c3B|ABAF E2 D2|A2 A2 cABc|dBAF "D"D2 DF|
"A"A2 A2 c3B|ABAF E2 FE|"D"D2 ED "G"FDGD|1"A"FDEF "D"D4:|2"A"FDEF "D"D2||
\end{abc}

\vskip 0.5in
\noindent\makebox[\textwidth][c]{%
\begin{minipage}{6.5in}
  \setlength{\parindent}{1.5em}
  
This tune is thought to have originated in the Upper South: Virginia,
West Virginia, and Tennessee. Each player had their own version of the
tune varying from 2 to 6 parts, with the earliest known printed
version of Forked Deer found in George P. Knauff's \textit{Virginia
Reels}, printed in 1839. The title's origin is ambiguous, with some
claiming it was named for the Forked Deer River in West Tennessee
(pronounced ``Fork-éd''), or as a corruption of ``Fauquier Deer'', the
name of a county in northern Virginia. It could also be a corruption
of ``Forked Air'', meaning a crooked tune, and some early sources do
have it as ``Forkadair'' or ``Forkedair Jig''. One version, from the
Hamblen Collection\footnote{http://chriswig.com/hamblen.html}, even
has the title ``Forked Ear'', although that must come from a
mis-hearing of the title. Perhaps more entertainingly, ``forked''
could be from an Ozarks region colloquial term for ``horny'' or
``sexually excited''.

``Forked Deer'' is one of the tunes named in Richard L. Dawson's 1885
poem \textit{The Hoosier Fiddle}\footnote{https://mudcat.org/thread.cfm?threadid=166702},
recalling the dances of his youth. Earlier sources disagree on
which part to begin the tune on, but by the time the tune was
established as a bluegrass standard it seems to have settled on being
a two-part tune with the high part played first, as set here.

\end{minipage}
}


\end{document}
